% !PRINTED 126
\weberpage{126}
%<<CONT

von Null verschieden; es mögen also $f'(\alpha),f''(\alpha),\ldots,f^{(m-1)}(\alpha)$ verschwinden, $f^{(m)}(\alpha)$ nicht verschwinden. Wir haben dann nach \S.\,13
\begin{equation}
 f(\alpha+h)=f(\alpha)+\frac{h^m}{\Pi(m)}f^{(m)}(\alpha)
 +\frac{h^{m+1}}{\Pi(m+1)}f^{(m+1)}(\alpha)+\cdots .
\tag{12}
\end{equation}

Wir wählen, was nach \S.\,31, V stets möglich ist, eine positive Zahl $\varepsilon$ so, dass
\begin{equation}
\left|\frac{h}{m+1}\frac{f^{(m+1)}(\alpha)}{f^{(m)}(\alpha)}
+\frac{h^2}{(m+1)(m+2)}\frac{f^{(m+2)}(\alpha)}{f^{(m)}(\alpha)}+\cdots\right|<1
\tag{13}
\end{equation}
ist, sobald
\begin{equation}
 |h|<\varepsilon,
\tag{14}
\end{equation}
und einen positiven echten Bruch $\delta$ so, dass
\begin{equation}
 \delta |f(\alpha)|<\left|\frac{\varepsilon^m f^{(m)}(\alpha)}{\Pi(m)}\right|,
\tag{15}
\end{equation}
was, wenn $f(\alpha)$ nicht verschwindet, gleichfalls möglich ist. Dann bestimmen wir (auf Grund von \S.\,32) $h$ aus der Gleichung
\begin{equation}
 \frac{h^m f^{(m)}(\alpha)}{\Pi(m)}=-\delta f(\alpha),
\tag{16}
\end{equation}
woraus nach (15)
\[
 |h|<\varepsilon
\]
folgt. Aus (12) ergiebt sich jetzt, wenn man für $h^m$ den Werth aus (16) setzt,
\[
 f(\alpha+h)=(1-\delta)f(\alpha)
 -\delta f(\alpha)\left(\frac{h}{m+1}\frac{f^{(m+1)}(\alpha)}{f^{(m)}(\alpha)}
 +\frac{h^2}{(m+1)(m+2)}\frac{f^{(m+2)}(\alpha)}{f^{(m)}(\alpha)}+\cdots\right)
\]
also
\begin{align}
 |f(\alpha+h)|\le{}& |f(\alpha)|(1-\delta)\notag\\
 &+\delta |f(\alpha)|\left|\frac{h}{m+1}\frac{f^{(m+1)}(\alpha)}{f^{(m)}(\alpha)}
 +\frac{h^2}{(m+1)(m+2)}\frac{f^{(m+2)}(\alpha)}{f^{(m)}(\alpha)}+\cdots\right|,
\tag{17}
\end{align}
und wegen (13)
\begin{equation}
 |f(\alpha+h)|<|f(\alpha)|\tag{18}
\end{equation}
\noindent\hfill w. z. b. w.

Damit ist also der Beweis des Fundamentalsatzes beendigt.
